\chapter*{Abstract}
\thispagestyle{specialheader}
\markboth{}{}
\addcontentsline{toc}{chapter}{Abstract}
This thesis presents a multi-objective optimisation of a compact
pressurised water reactor core fuelled with low-enriched \ce{UO2}, based
on surrogate models and the OpenMC Monte Carlo code. The objective was to find
core configurations that meet the requirements known from open sources
for the land-based prototype of the Brazilian naval propulsion
programme, and to quantify the trade-off
between cycle length, radial peaking and the soluble boron needed for
criticality.

The prototype is rated at 48\,MW thermal and its core data are not
public. The model was built from the few facts in the open record, namely
the reactor type, the power, the fuel class, the vessel envelope and a
five-year refuelling interval, and from declared assumptions. The enrichment, the gadolinia loading and the reflector
thickness were varied on a single-batch core of 32 assemblies. Because one
depletion evaluation takes around 20 minutes and carries statistical
noise, a Gaussian-process surrogate model trained on the Monte Carlo
evaluations was searched by the NSGA-II multi-objective genetic algorithm
in an active-learning loop over nine campaigns and
636 evaluations, and the candidates were confirmed on a
three-dimensional core model.

Campaigns 1 to 8 maximised the cycle length and minimised the radial
peaking factor. The Campaign 8 champion, C8-47, reaches 1941 effective
full power days at a peaking factor of 1.510, but every design on that
front exceeded the five-year requirement, and the longer cycles were
bought with excess reactivity that only deeply inserted control banks
could hold. Campaign 9 therefore imposed the cycle length as a
constraint at 1826 effective full power days and minimised the peaking
factor and the critical boron concentration at the most reactive state.
Its front holds five designs at 4.2 to 4.5\,wt\% enrichment, with peaking
factors of 1.490 to 1.545 and boron demands of 1406 to 1908\,ppm, and the
boron demand falls as the gadolinia loading rises.

These results hold for a two-dimensional beginning-of-life model in the
loop with fixed temperatures and constant boron in the depletion, and
neighbouring designs lie within the Monte Carlo noise on the peaking axis.
On the three-dimensional core model C9-47, C9-34 and C9-40 stay on the
Campaign 9 front, and all five designs are held subcritical by the first
two regulating banks.
